% Figure: two op-amp configurations, inverting (left) and non-inverting
% (right), drawn with a plain TikZ triangle for the amplifier symbol.
\begin{figure}[htbp]
\centering
\begin{tikzpicture}[
  line width=0.8pt,
  res/.style={draw, fill=white, minimum width=11mm, minimum height=4mm, font=\scriptsize},
  font=\small,
]
% --- Inverting amplifier (left) ---
\begin{scope}[xshift=0cm]
  % op-amp triangle
  \draw[engblue] (1.6,0.6) -- (1.6,2.4) -- (3.4,1.5) -- cycle;
  \node[font=\scriptsize, anchor=west] at (1.7,2.0) {$-$};
  \node[font=\scriptsize, anchor=west] at (1.7,1.0) {$+$};
  % input resistor Rin
  \node[res] (Rin) at (0.4,2.0) {$R_{\text{in}}$};
  \draw[engblue] (-0.8,2.0) -- (Rin) -- (1.6,2.0);
  \node[font=\scriptsize, anchor=east] at (-0.85,2.0) {$v_{\text{in}}$};
  % feedback resistor Rf
  \node[res] (Rf) at (2.5,3.2) {$R_f$};
  \draw[engblue] (1.6,2.0) -- (1.6,3.2) -- (Rf) -- (3.9,3.2) -- (3.9,1.5);
  % output
  \draw[engblue] (3.4,1.5) -- (4.4,1.5);
  \node[font=\scriptsize, anchor=west] at (4.45,1.5) {$v_{\text{out}}$};
  % non-inverting input grounded
  \draw[engblue] (1.6,1.0) -- (1.0,1.0) -- (1.0,0.4);
  \node[font=\scriptsize] at (1.0,0.15) {gnd};
  \node[font=\small] at (1.8,-0.4) {inverting, gain $-R_f/R_{\text{in}}$};
\end{scope}
% --- Non-inverting amplifier (right) ---
\begin{scope}[xshift=7cm]
  \draw[engblue] (1.6,0.6) -- (1.6,2.4) -- (3.4,1.5) -- cycle;
  \node[font=\scriptsize, anchor=west] at (1.7,2.0) {$-$};
  \node[font=\scriptsize, anchor=west] at (1.7,1.0) {$+$};
  % input to non-inverting terminal
  \draw[engblue] (0.4,1.0) -- (1.6,1.0);
  \node[font=\scriptsize, anchor=east] at (0.35,1.0) {$v_{\text{in}}$};
  % feedback divider Rf and Rg
  \node[res] (Rf2) at (2.5,3.2) {$R_f$};
  \draw[engblue] (1.6,2.0) -- (1.6,3.2) -- (Rf2) -- (3.9,3.2) -- (3.9,1.5);
  \node[res, rotate=90] (Rg) at (1.6,2.9) {};
  \node[font=\scriptsize, anchor=east] at (1.35,2.9) {$R_g$};
  \draw[engblue] (1.6,2.0) -- (Rg);
  \draw[engblue] (Rg) -- (1.6,3.5) -- (1.0,3.5) -- (1.0,3.9);
  \node[font=\scriptsize] at (1.0,4.15) {gnd};
  % output
  \draw[engblue] (3.4,1.5) -- (4.4,1.5);
  \node[font=\scriptsize, anchor=west] at (4.45,1.5) {$v_{\text{out}}$};
  \node[font=\small] at (1.8,-0.4) {non-inverting, gain $1+R_f/R_g$};
\end{scope}
\end{tikzpicture}
\caption[Inverting and non-inverting op-amp configurations]{The two basic
negative-feedback op-amp amplifiers. Both follow from the virtual-short
constraints of the ideal op-amp: the two input terminals sit at the same
voltage and draw no current. The inverting configuration has gain
$-R_f/R_{\text{in}}$ and input impedance $R_{\text{in}}$; the non-inverting
configuration has gain $1 + R_f/R_g$ and very high input impedance. The
choice between them turns on the sign of the gain and on the source impedance
the amplifier must present to its driver.}
\label{fig:vol04ch09:opamp-configs}
\end{figure}
